\chapter{Results}\label{chapter:results}

Having introduced all the necessary components we can start looking at the results.
What follows is a rough outline of the results chapter with reference to the previous chapters and section in which we introduced the necessary concepts.

In the results chapter we look at the thresholds (see section \ref{ssec:threshold_theorem}) of different configurations.
All of those configurations are encoded in the surface code (see section \ref{sec:kitaev_surface_code}) with varying code distance $d$.
We change the applied error model (see section \ref{ssec:error_models}) and the syndrome extraction circuit (see section \ref{ssec:syndrome_extraction}),
with the goal of characterizing the Steane-type error correction circuit (see section \ref{sssec:steane_type_error_correction_intro}). 

In the first section (section \ref{sec:results_code_capacity}) we will test our implementation of STEC by comparing the observed threshold of the simulation under the code capacity noise model, 
with analytically predicted thresholds, to show that everything works as expected.
Here we will also highlight the performance and problems with ML decoding, namely the numerical underflow issue, further discussed in the appendix \ref{chapter:numerical_underflow}, 
and we will go through an example of the data collapse method of finite size scaling analysis, as introduced in section \ref{eq:data_collapse_method}, applied to our simulated data.

This is followed by the characterization of STEC under circuit level noise, in section \ref{sec:results_circuit_lvl}.
Here we use the conventional syndrome extraction circuit, as described in section \ref{sssec:standard_syndrome_extraction_circuit}, as a reference to which we compare STEC.
Looking at the threshold of a single QEC cycle utilizing STEC syndrome extraction, in section \ref{ssec:results_circ_single_cycle}, 
we can highlight the observed difference in ordering and show the influence that the choice of decoding algorithm has on the results. 
We will discuss separately how the decoding algorithm of the residual error, which is required as described in section \ref{sec:fault_tolerance}, influences the results in subsection \ref{sssec:ml_ft_decoding}.

We can then further characterize STEC by simulating and analyzing multiple consecutive independent QEC cycles, as shown in section \ref{ssec:results_circ_multiple_cycles}.
This analysis shows the relevance of the incoming error and points towards a stable threshold, that is reached after multiple rounds.
Comparing the STEC approach to the conventional syndrome extraction, or at least to the upper bound of it, due to our simplifying assumptions as described in \ref{ssec:conv_se_circ_lvl_noise_adaption}, 
will give us an estimate that the threshold of STEC is higher thant that of the conventional approach by a factor of approximately two.

The final conclusion of those results will then be drawn in chapter \ref{chapter:conclusion}.

\section{Code Capacity}\label{sec:results_code_capacity}

We start our result chapter by simulating the STEC extraction circuit with a code capacity noise model, 
to get thresholds which we can compare to predicted analytical threshold for the ML \cite{topological_quantum_memory} and MWPM \cite{efficient_decoding_algorithm,MWPM_threshold_paper} case.

We use in this case four different configuration, consisting of two possible choices of observable, and corresponding error model, and two choice of decoding algorithms.
In this simulation we always use the STEC extraction circuit.
For the observable we either measure the logical $\bar{Z}$ observable, in combination with a code capacity model implementing bit-flip errors ($X$-errors) after the initialization,
or we measure the $\bar{X}$ observable, while simulating a circuit with phase-flip errors ($Z$-errors) as the noise source.
Regarding the decoder algorithm we use the two decoding algorithm introduced in section \ref{sec:decoding_problem_general}, either MWPM or ML decoding.
Note that for this case the ML decoder is exact, as it assumes either bit-flip or phase-flip noise, which is the case in this example.

For each of those measurement series we perform the analysis, described in the threshold chapter \ref{ssec:threshold_theorem}, to determine the threshold.
To show the intermediate results and highlight some aspects of this approach we have discussion and describe the process and its intermediate results in section \ref{ssec:results_example_process}.



% Discuss results first before going into details about how those are obtained?!
% good agreement 
% ml might show finite size effects, might want to switch to distance 7 (show 3,5,7) 


\subsection{Example Process and Intermediate Results}\label{ssec:results_example_process}

In the following we will show and discuss the analysis and intermediate results in the process to determining the threshold of the configuration 
where we observed logical $\bar{Z}$ observable and decode the syndrome using the efficient ML decoder.

We use the \texttt{stim} library to construct and simulate the quantum circuit, with the designated noise model, to obtain the syndrome and observable measurement, the logical $\bar{Z}$ in this case. 
The syndromes are then decoded using the decoder of choice, adapted to this noise model and circuit.
Utilizing Pauli frame tracking, introduced in section \ref{sec:pauli_frame_tracking}, we can check if the measurement agrees with the predicted recovery or if the recovery would have been unsuccessful.
Based on those results we can determine the logical error rate $p_{log}$.

We repeat this simulation multiple times, while varying the code distance $d$ of the surface code of the simulated quantum circuit, and the physical noise rate $p_{phy}$ of the noise model.
Each of those simulations is itself repeat a number of times which we will refer to as the number of shots $n_{shots}$, where each shots results in a logical observable that need to be predicted.

We can then plot the logical error rate $p_{log}$ against the physical error rate $p_{phy}$ for the multiple code distance $d$ we simulated. 
As explained in section \ref{sec:kitaev_surface_code}, we limit ourselves to odd distance codes and we furthermore start our analysis at distance $d=5$ to avoid some finite size effects. 
We will discuss those finite size effects in more detail towards the end of this chapter.
For the physical error rate $p_{phy}$ we limited by physics to the interval $0 < p_{phy}<1 $.
To use reasonable physical error rates $p_{phy,i}$ for the noise model we use physical error rates which are equally spaced on a logarithmic scale around an expected threshold $\tilde{p}_{th}$
\begin{equation}
    p_{phy,i} = \tilde{p}_{th} \cdot 10^{k\cdot i}, \ 0 < k < 1, \ i \in \{-N,...,N\}.
\end{equation}
We will further limit the interval in a later step of the analysis. 
% name scaling convention for physical error rate, logspace, equaly distacne .... so on and so on 
The results can be seen in figure \ref{fig:results_code_capcity_Z_ml_p_log_p_phy} for the logical $\bar{Z}$ observable, under circuit level noise and using the efficient ML decoder.

We can spot the approximate location of the threshold as the point where the logical error rates of different code distances cross, as described in section \ref{ssec:threshold_theorem}. 
The asymptotic behavior described by equation \ref{eq:p_log_asymptotic_behavior}, can also be identified toward low physical error rates.


\begin{figure}[t]
    \centering
    \includegraphics[width=0.6\linewidth]{graphics/circ_ml_z_log_error_rate_complete_rejection_rate.pdf}
    \caption{
        \textcolor{red}{update this image to be the correct image!!!}
        }
    \label{fig:results_code_capcity_Z_ml_p_log_p_phy}
\end{figure}

Given that we here used the ML decoder, we encounter the numerical underflow problem in this implementation, as shown by the faint blue line in figure \ref{fig:results_code_capcity_Z_ml_p_log_p_phy}.
The faint blue line shows the rejection rate and we reject those shots in which we detected the numerical underflow problem during the decoding process.
Here we can state that in the interval in which we are interest in, the interval around the thresholds, the numerical rejection rate is low enough to be ignored.
This problem and possible solutions are further discussed in the appendix in section \ref{chapter:numerical_underflow}. 
Note that the MWPM decoder does not suffer from the same problem. 

To correctly determine the threshold we continue our analysis with the data collapse method.
For this we limit the interval of interest to the one shown in figure \ref{fig:results_code_capcity_Z_ml_p_log_p_phy_small_window}.
This interval is chosen to be centered around the observed threshold and limited to be close to the phase transition for the finite size scaling analysis.

\begin{figure}[t]
    \centering
    \includegraphics[width=0.6\linewidth]{graphics/circ_ml_z_log_error_rate_complete_rejection_rate.pdf}
    \caption{
        \textcolor{red}{update this image to be the correct image!!! should be cut out of plog against pphy of the interval used for fssa}
        }
    \label{fig:results_code_capcity_Z_ml_p_log_p_phy_small_window}
\end{figure}

The data collapse method, described in section \ref{eq:data_collapse_method}, collapse the data points of the chosen interval, shown in figure \ref{fig:results_code_capcity_Z_ml_p_log_p_phy_small_window},
onto a single curve and in this process the critical exponent and the threshold are determined. 
Looking at the collapsed data point on a curve, as shown in figure \ref{fig:results_code_capcity_Z_ml_p_log_p_phy_collapsed_data}, is an indication if this method was successful.
We therefore show the curves of all of those fitting used in this thesis in the appendix in section \ref{appendix:data_collapse_figures}. 
For this case the data collapse method worked and we were able to correctly determine the threshold.

\begin{figure}[t]
    \centering
    \includegraphics[width=0.6\linewidth]{graphics/circ_ml_z_log_error_rate_complete_rejection_rate.pdf}
    \caption{
        \textcolor{red}{update this image to be the correct image!!! should be collapsed curve}
        }
    \label{fig:results_code_capcity_Z_ml_p_log_p_phy_collapsed_data}
\end{figure}

This method can fail if we use a wrong initial estimate for the threshold.
Including the distance $d=3$ also leads to a worse behavior, as the finite size effects become stronger and our assumption behind the data collapse method become less valid.
\textcolor{red}{show example for low distances?}


\section{Circuit Level}\label{sec:results_circuit_lvl}

% shortly introduce circuit level as the design where the syndrome extraction circuit difference becomes relevant
% motivate again why circuit level and refer back to previous section about circuit level encoding

% shortly discuss the numerical underflow issue of the decoder then point to the appendix for further reference and discussion

\subsection{Single QEC Cycle}\label{ssec:results_circ_single_cycle}
% repeat the setup as described in the previous chapter 

% highlight use of mwpm perfect matching as a residual weight decoding approach

% Show single round circuit level noise behaviour 

% point out the differences:
% difference in decoding, once again
% difference in ordering -> hard to judge apriori which should be higher

% lead to influence of incomming error -> motivated repeated syndrome extraction and decoding 

% compare to conventional syndrome extraction circuit,
% comment on the different amount of ressources required 
% name again the simplfying assumption of conventional decoding!

\subsubsection{ML Decoding Residual Error}\label{sssec:ml_ft_decoding}

% discuss decoding of residual error using an ml decoder
% error rate approximated using the physical error rate, not exact just testing

\subsection{Multiple independent QEC Cycles}\label{ssec:results_circ_multiple_cycles}

% talk about implementation 

